% Vietnamese translation for OI Public Library
% Source-File: IOI中国国家候选队论文/2004/鬲融.ppt
\documentclass[11pt]{article}
\usepackage{oipl-vn}

\newcommand{\Oh}{\mathcal{O}}

\title{Tư tưởng vét cạn đặc biệt\\{\large Ghi chú theo slide}}
\author{Ge Rong}
\date{IOI 2004}

\begin{document}
\maketitle

\origtitle{Special enumeration thinking}
\sourcepath{IOI 2004 / Ge Rong.ppt}

\section{Vét cạn không phải lúc nào cũng thô}

Bài trình bày nhấn mạnh rằng vét cạn vẫn hữu ích nếu không gian được mã hóa
đúng và đủ nhỏ. Thay vì duyệt dữ liệu đầu vào, ta duyệt một tập trạng thái
tham số đại diện cho toàn bộ hành vi có thể xảy ra.

\section{Smart Typist}

Slide nhắc bài \emph{Smart Typist} của NOI 2001 với các thao tác
\texttt{Up/Down}, \texttt{Left/Right}, \texttt{Swap0/Swap1}. Chuỗi ví dụ
\texttt{123456} và \texttt{633451} cho thấy bài toán có thể nhìn như hoán vị
trên \(6\) ký hiệu. Dùng băm và A* trên \(6!=720\) trạng thái, mỗi trạng thái
có khoảng \(10\) chuyển tiếp, cho chi phí hằng số theo nghĩa thực tế của bài.

\section{Logic Island và Lock}

Một ví dụ khác là \emph{Logic Island}: các mệnh đề dạng ``tôi là hoặc không
là một loại người'', ``người \(X\) là hoặc không là một loại người'', và
``đang là ngày hoặc đêm''. Slide nhắc không gian khoảng \(3^5\cdot 2=486\)
khả năng, nên có thể duyệt mọi cách gán rồi kiểm tra tính nhất quán trong
\(\Oh(s)\) theo số câu.

Bài \emph{Lock} của NOI 2003 xuất hiện với ký hiệu \(T(i,j)\) và bước cải
thiện từ \(\Oh(n^5)\) xuống \(\Oh(n^3\log n)\). Đây là ví dụ về việc vét
cạn tham số biên rồi dùng thuật toán hiệu quả cho phần còn lại.

\section{Ghi nhớ}

Vét cạn đặc biệt đòi hỏi tìm không gian nhỏ nhưng đủ biểu đạt. Nếu trạng thái
đại diện đúng, thuật toán có thể vừa đơn giản vừa mạnh.

\end{document}
